\section{Introduction}
\label{sec:intro}

Quantized LLM agents are usually compared with their full-precision reference by average task success, and in closed-loop benchmarks such as \tautwo{} \citep{barres2025tau2} that comparison can look reassuring. On our checkpoints, quantized-minus-BF16 success ranges from $-0.2$ to $+4.2$\pp{} across four W4 cells, with new-task (crossed) intervals about $\pm$6--7\pp{} wide (Appendix~\ref{app:e2e}). % src: analysis/insights/iter_4/out/success_ni_ALIGNED_20260930T213009Z.md
Yet the quantized agent (Q) still takes actions that its BF16 reference (B) would not. If B is available as a fallback, as in dynamic precision routing \citep{li2026dynamicmp}, every such state poses a decision: \textbf{upgrade precision for this step only, from the next step on, or both?} In a closed-loop agent, the value of such an upgrade is jointly determined by the action it changes and by the execution that follows.

\paragraph{Judging the two actions is not enough.} Consider a state from our data (q4 AWQ-W4, \tautwo{} airline task 31, nonce block f13, decision 1 of 12). Q looks up a reservation id that does not exist and receives ``not found''; B looks up the user's profile. B's action plainly looks better. We froze this state and ran all four combinations of \{B's or Q's action\} $\times$ \{B or Q continuing\}. Only one succeeds: B's action \emph{followed by} B. If B supplies only this step and Q continues, Q searches for flights, writes a change and fails. If Q's failed lookup stands and B continues, B also fails. % src: analysis/fork/out/ghf_20261001T211519Z.md (§4 Case 5)
At this state, the better-looking action has task value under one continuation and none under the other; judging the action alone cannot reveal this.

\paragraph{A crossed same-state intervention.} At a frozen state whose no-op replay is bit-exact, we cross the source of the current action with the controller of the remaining episode (Figure~\ref{fig:fork}). This turns ``this step'' ($G$), ``from the next step'' ($H$) and ``both'' ($F$) into directly comparable interventions on the same states (\S\ref{sec:method}). The $2\times2$ crossing is a standard factorial design. We study Qwen3-4B and Qwen3-8B (simulated weight and KV-cache quantization, BF16 compute) in \tautwo{} with one greedy user simulator, BFCL multi-turn and ALFWorld, forking at the first decision where the two models' actions differ in tool structure (text-only differences are skipped) and at ordinary states drawn independently of any outcome. % nocheck: design label and model sizes

\paragraph{What we find.}
\begin{itemize}[leftmargin=*,itemsep=1pt,topsep=2pt]
\item \textbf{The value of an action depends on the controller that continues} (\S\ref{sec:r2}). In 26.9\% of development W4 divergent states, the same action swap has a different outcome effect under B's continuation than under Q's. The two effects agree more than a marginal-preserving pairing control expects (38.0\%, $\kappa=0.29$); strict reversals are rare. Post hoc, it also appears on BFCL (4.7--17.1\% per W4 cell) and q8 ALFWorld (32.8--40.0\%), which have no LLM user, and on the telecom holdout (20.7\%). Without a same-precision perturbation control, these shares include the continuation's sensitivity to any change of controller. Cross-cell comparisons are descriptive, and no $G-H$ ordering is established as transferable. % src: analysis/fork/out/ghf_20261001T211519Z.md (§2); analysis/fork/out/agreement_control_20261002T024240Z.md (§2); analysis/fork/out/ext_r2_20261002T024433Z.md (§3); analysis/fork/out/holdout_r2_posthoc_20261002T024414Z.md (§3 pool)
\item \textbf{Local precision upgrades redistribute success and failure} (\S\ref{sec:r1}, \S\ref{sec:r3}). Replacing a W4 agent's first tool-structure-divergent action rescues or damages 15.5--20.9\% of task$\times$block units, while the net effect is $-1.9$ to $+2.9$\pp{}. On 34 pre-registered unseen telecom tasks, the swap again rescues and damages, task fragility repeats, and the estimated effect is $-4.6$\pp{} (90\% CI $-7.2$ to $-2.0$; the frozen primary reading is inconclusive). A local precision patch therefore needs a direct evaluation of its execution consequences. % src: analysis/fork/out/dev_all_20261001T070326Z.md; analysis/fork/out/holdout_confirm_20261002T012353Z.md
\item \textbf{Finding sensitive decisions and choosing beneficial replacements are different learning targets} (\S\ref{sec:r4}). A frozen Q-uncertainty trigger captures more outcome-changing decisions than random selection at two of three rates on later nonce blocks, but has not been shown to select net-beneficial replacements; a hindsight oracle with the same per-cell budget nets $+3.3$ to $+3.7$\pp{}. Mixed-precision selection thus has to predict rescues and damages jointly, not divergence alone. % src: analysis/fork/out/token_mech_20261002T003130Z.md (§3); analysis/fork/out/oracle_trigger_20261002T024423Z.md (§1a)
\end{itemize}

\paragraph{Contributions.} (1)~A crossed-intervention evaluation framework for precision upgrades in agents: same-state restoration with bit-exact replay acceptance, full-action interventions, paired estimands ($G$, $H$, $F$ and the $(\DQ,\DB)$ table) with explicit denominators, and inference that separates a fixed task pool from new tasks. (2)~Systematic findings over weight, KV-cache and stress quantization, two model sizes and three benchmarks: two-sided consequences, controller dependence across environments, text actions at ordinary states (labels not reproduced under a second simulator; \S\ref{sec:r3}) and diagnostics for precision selection. Withdrawn results are in Appendix~\ref{app:withdrawn}. % nocheck: list labels (1)-(2)

Evaluating a precision upgrade for an agent thus requires measuring the task value of the replacement under the execution that follows it, rather than the replaced action in isolation or either model's average success.

\begin{figure}[t]
\centering
\begin{tikzpicture}[ % nocheck: TikZ layout coordinates
  font=\footnotesize,
  box/.style={draw, rounded corners=2pt, align=center, inner sep=2.5pt},
  act/.style={draw, rounded corners=2pt, align=center, inner sep=2pt, fill=gray!12}, % nocheck: TikZ style
  br/.style={draw, align=center, inner sep=1.5pt, minimum width=1.3cm},
  arr/.style={-{Stealth[length=4pt]}, thin}]
\node[box, text width=7.0cm] (h) at (0,0) {frozen state $s_t$: history, environment and device DBs, user-simulator view, used step/error budget}; % nocheck: TikZ
\node[act] (aB) at (-1.9,-1.0) {current action $a_{\mathrm{B}}$}; % nocheck: TikZ
\node[act] (aQ) at (1.9,-1.0) {current action $a_{\mathrm{Q}}$}; % nocheck: TikZ
\draw[arr] (h.south -| aB) -- (aB); \draw[arr] (h.south -| aQ) -- (aQ);
\node[br] (bb) at (-2.85,-2.0) {BB\\[-1pt]{\scriptsize then B}}; % nocheck: TikZ
\node[br] (bq) at (-0.95,-2.0) {BQ\\[-1pt]{\scriptsize then Q}}; % nocheck: TikZ
\node[br] (qb) at (0.95,-2.0) {QB\\[-1pt]{\scriptsize then B}}; % nocheck: TikZ
\node[br] (qq) at (2.85,-2.0) {QQ\\[-1pt]{\scriptsize then Q}}; % nocheck: TikZ
\draw[arr] (aB) -- (bb); \draw[arr] (aB) -- (bq);
\draw[arr] (aQ) -- (qb); \draw[arr] (aQ) -- (qq);
\foreach \n/\l in {bb/BB,bq/BQ,qb/QB,qq/QQ} {\node[below=0.08cm of \n, font=\scriptsize] (y\n) {$Y_{\mathrm{\l}}$};} % nocheck: TikZ
\node[box, text width=7.0cm, fill=orange!8, anchor=north] (cat) at (0,-2.8) % nocheck: TikZ
  {$G=\YY{BQ}-\YY{QQ}$ (this step only),\; $H=\YY{QB}-\YY{QQ}$ (from the next step)\\[1pt]
   $F=\YY{BB}-\YY{QQ}$ (from here on);\quad $\DQ=G$,\; $\DB=\YY{BB}-\YY{QB}$};
\end{tikzpicture}
\caption{Crossed same-state intervention (schematic, not data). B = BF16 agent, Q = quantized agent. In branch XY, the first letter is the source of the current action and the second is the controller from the next step to the end. Every branch starts from the same frozen state with the original remaining budget. QQ re-executes Q's original episode, and its bit-exact reproduction is the no-op check. Definitions are in \S\ref{sec:estimands}.}
\label{fig:fork}
\end{figure}
